%!TEX root = ../template.tex %
% !TeX spellcheck = en_US
%%%%%%%%%%%%%%%%%%%%%%%%%%%%%%%%%%%%%%%%%%%%%%%%%%%%%%%%%%%%%%%%%%%%
%% chapter4.tex
%% NOVA thesis document file
%%
%% Chapter with lots of dummy text
%%%%%%%%%%%%%%%%%%%%%%%%%%%%%%%%%%%%%%%%%%%%%%%%%%%%%%%%%%%%%%%%%%%%

\typeout{NT FILE chapter7.tex}%

\chapter{Outlook}
\label{ch:Outlook}

Brief chapter on what else could be done with the code.

\section{Quasi-2D approach: the multilayer case}
\label{sec:quasi_2d_approach}

Here I refer to the interaction matrix elements in the mixed representation, for the multilayer case.
Probably put that as appendix, and join this part with the real space approach one, both can deal with the multilayer case.

\section{Real space approach to electrostatic screening}
\label{sec:real_space_approach}

Probably put this in appendix.

\section{Dielectric screening in 2D metals}
\label{sec:dielectric_screening_2d_metals}

An obvious extension of the code is to include metals, for example graphene.
One could compare the numerical results with those the analytical ones, plenty of literature for those.
Test the role of local-field effects in graphene, as we did for insulators.

\section{Non-local excitonic conductivity beyond the dipole approximation}
\label{sec:non_local_exciton_conductivity}

Having the plane-wave matrix elements for any $\mathbf{Q}$, it could eventually be implemented in the code.
This would allow also computing the radiative decay of excitons in extended systems beyond the dipole approximation, never done in the literature.
The study of the nonlocal transversal conductivity beyond the dipole approximation remains unexplored.
Calculation of the matrix elements $\langle n,\bk |\ee^{\ii \bq \vdot \br} \bp | n', \bk'\rangle$ by inserting resolution of identity.